
% Authoritative LaTeX manuscript; Markdown is a generated compatibility projection.

\textbf{Problem.} A model becomes dangerous when engineers and agents continue to trust it after it no longer corresponds to the system it represents.

\textbf{Move.} Treat correspondence as an invariant. Check both that reality still satisfies modeled claims and that consequential implementation elements have not appeared outside the model\textquotesingle s declared scope.

\autoref{fig-move03} shows the two directions.

\begin{figure}[htbp]
\centering
\includegraphics[width=\linewidth]{book/latex/assets/c3-bidirectional-correspondence-82672ca26c.pdf}
\caption{\emph{Correspondence runs both ways.} A two-way loop joins MODEL and REALITY. One arrow asks whether reality still satisfies the model, catching a wrong modeled fact; the other asks whether something important appeared outside the model, catching unmodeled reality.}
\label{fig-move03}
\end{figure}

\textbf{Example --- Architecture drift.} A component-and-zone model declares where code belongs and which dependencies it may take. A gate compares the declared architecture against the imports the implementation actually makes. A change that violates the modeled structure fails the check at commit time.

\textbf{Example --- Governance coverage.} The same correspondence idea runs the other way. An orphan-coverage walk starts from code and asks what model or control governs it. Uncovered code is not automatically wrong; it shows where implementation currently sits outside the modeled governance estate and requires an explicit decision.

\textbf{Related mechanisms:} Component \& zone model · Drift \& parity gates · Cross-source coherence lints · Orphan-coverage metric · Symbol-anchored traceability.
